\section{Preserving admission when the period is moved}

% BEGIN W75A1
\paragraph{W75A1. The finite level and norm after a bounded-denominator move.}
Fix $c,R>0$, an even row count $n$, and a rational split unitary
matrix $g$ of degree $d$ with $cg$ integral. Its inverse has the
same denominator bound by the unitary identity. Let $L$ be an
original definite self-dual lattice and $A=L^d$ the additive lattice
of test matrices with columns in $L$. Let $D=|\operatorname{disc} E|$,
$T=2cR$ and $R'=DT$. Suppose $f$ is supported on $R^{-1}A$ and
invariant by $RA$. Then $f_g=\omega(g)f$ has support in
$(R')^{-1}A$, has period $R'A$, and satisfies
\[
              \|f_g\|_\infty
              \le (RT)^{[E:\mathbb Q]nd/2}\|f\|_\infty .
\]
All norms in this assertion use the original finite Weil action.

Here it is convenient to read a rational Schwartz--Bruhat test on
the finite adeles. This amounts to extending its values on the
finite quotient of a denominator lattice by a period lattice.
Completion gives the same finite quotient, by the Chinese remainder
theorem. We use hats for the resulting compact open lattices. The
native Heisenberg operators extend to these completions by local
constancy on every test. The covariance identity for rational
displacements extends by density in the same way.

Displacements with both translation and modulation in $2R\widehat A$
fix $f$. In the phase their half-integral cross term and the
modulation against an argument of the support both have integral
trace, by the Hermitian self-duality of $L$. Since both blocks of
a displacement in $T\widehat A\oplus T\widehat A$ multiplied by
$g$ are in $2R\widehat A$, covariance makes that displacement fix
$f_g$. Pure translations give its period $T\widehat A$.
Pure modulations force its support into the trace dual of
$T\widehat A$. The inverse different is the trace dual of
$\mathcal O_E$. Testing all integral multiples of each Hermitian
pairing therefore places this support in
$T^{-1}\mathfrak D_E^{-1}\widehat A$. Indeed locally the dual
basis for the Hermitian pairing identifies this lattice with the
inverse different in each coordinate. The integer $D$ clears the
inverse different, giving the required support and period.
On each coset of $T\widehat A$, the square of the point value times
its volume is at most $\|f_g\|_2^2=\|f\|_2^2$. The latter is at
most $\operatorname{vol}(R^{-1}\widehat A)\|f\|_\infty^2$.
Their relative index is exactly $(RT)^{[E:\mathbb Q]nd}$.
This proves the estimate for every point and lattice.
% END W75A1

% BEGIN W75A2
\paragraph{W75A2. The locality field survives the same Weil move.}
If $\phi$ is admitted at $(M,R)$, the family
$L\mapsto\omega(g)\phi_L$ is admitted at $(M,R')$ for the $R'$
of W75A1. In particular the bound on the level and the supremum
bound do not depend on the chosen definite lattice, test or move.

Positivity, support and periods follow from W75A1. A rational unitary
matrix on the definite rows commutes with the column Weil action:
reindexing the defining theta transformation on the rows changes
its two elliptic variables contragrediently and leaves its scalar
factor invariant. Fourier uniqueness on all elliptic arguments
identifies the transformed tests, as in P75H3. This transports the
given equivariance of $\phi$.

For locality suppose $L,L'$ agree at $MR'$. Let $S$ be the finite
set of rational primes where their completions differ. It is
disjoint from the primes dividing $2cRDM$. Set
$B=\prod_{p\in S}A_p$, $B'=\prod_{p\in S}A'_p$, and write
finite adelic arguments as $(y_S,y^S)$. On these factors both $R$
and $R'$ are units. Support and periods force
\[
 \phi_L=\mathbf1_B\otimes h,\qquad
 \phi_{L'}=\mathbf1_{B'}\otimes h
\]
with a common function $h$ away from $S$. Indeed evaluate at
$y_S=0$. Outside $S$ the lattices coincide. On their allowed support,
admitted locality at $RM$ gives equality: choose a rational
argument in the same open neighborhood, with $S$-part in
$B\cap B'$ and with its other components in the desired coset of
the common period lattice. Such an argument exists by the same
finite-quotient approximation as in W75A1. Outside that support
both functions vanish. If $S$ is empty, this already gives equality.

Average the phase-space displacements over $B'\oplus B'$ with
probability Haar measure, acting on $S$ alone; call the operator
$P'$. At these odd unramified primes $B'$ is trace self-dual, by
Hermitian self-duality and triviality of the local different.
The half cross term of a displacement in this lattice is integral,
so averaging modulations is exactly multiplication by
$\mathbf1_{B'}$. Averaging translations then averages over $B'$.
The two averages commute. Consequently
\[
 P'\phi_L=a\phi_{L'},\qquad
        a=\frac{\operatorname{vol}(B\cap B')}{\operatorname{vol}(B')}>0.
\]
The same formula for the projector applies to any outside function
in the place of $h$. Because $g$ and $g^{-1}$ are integral at $S$,
right multiplication permutes $B'\oplus B'$ and preserves its
probability measure. Finite Heisenberg covariance therefore makes
$P'$ commute with $\omega(g)$. This reasoning can equally be
performed on a finite quotient on each smooth test.

Apply W75A1 to the two transformed tests. On $S$ their factors
must again be $\mathbf1_B$ and $\mathbf1_{B'}$ respectively,
times possibly different outside functions. The equality
$P'\omega(g)\phi_L=a\,\omega(g)\phi_{L'}$ and the projector
calculation show that those two functions coincide. On an argument
with $R'y$ integral in both lattices, its $S$-part belongs to their
intersection, so the transformed values agree. This is exactly
the remaining clause of admission, including changes of lattice
away from the chosen level.
% END W75A2
